\section{Universal coefficients for successive cycles}
\label{sec:universal}

The power construction applies to commutative algebras defined at the
individual primes. Our successive Koszul algebras, however, need only
exist in the rational mapping category. We connect these settings by
constructing a universal algebra for the coefficients of finitely many
cycles. Its generators have positive formal mass. This permits both a
rational comparison at fixed weights and a calculation of the same
primewise models at weights that increase with the prime.

Fix $0\leq a<n$, and suppose that cycles $b_j$, for $0\leq j\leq a$,
have been chosen in the successive strict extensions
$A[e_0,\ldots,e_{j-1}]$ of Section~\ref{sec:foundations}.
Here $|e_j|_{\mathrm{chain}}=1$, $\wt(e_j)=w_j$, and
$\partial e_j=b_j$. In particular, $b_0=\varpi u_0$.
We write $[j)=\{0,\ldots,j-1\}$ and take every exterior monomial
$e_S$ in increasing order of its indices. All signs below are the
signs of the homological grading; the formal lattice contributes no
additional sign.

\subsection{The coefficient algebra and its filtration}

\begin{definition}\label{universal:coefficients}
The graded algebra underlying $R^a$ is the free graded-commutative
$D$-algebra on variables
\[
 z_{j,S},\qquad 0\leq j\leq a,\quad S\subseteq[j),
 \qquad |z_{j,S}|_{\mathrm{chain}}=-|S|,
 \qquad \wt(z_{j,S})=\gamma_{j,S}
       :=w_j-\sum_{h\in S}w_h.
\]
Set
\[
 b_j^u=\sum_{S\subseteq[j)}z_{j,S}e_S,
 \qquad \partial e_j=b_j^u.
\]
The differential on $R^a$ is determined by
$\partial b_j^u=0$, with coefficients compared in the exterior basis.
Define
\[
 K_-^a=R^a[e_j:j<a],\qquad K^a=R^a[e_j:j\leq a].
\]
The superscript $u$ indicates these universal cycles.
\end{definition}

For example, when $n\geq2$ the stage $a=1$ has universal cycles
\[
 b_0^u=z_{0,\emptyset},\qquad
 b_1^u=z_{1,\emptyset}+z_{1,\{0\}}e_0.
\]
The coefficient $z_{1,\{0\}}$ has chain degree $-1$, whereas the
two empty-subset coefficients have chain degree zero. The cycle
equations require
\[
 \partial z_{0,\emptyset}=\partial z_{1,\{0\}}=0,
 \qquad
 \partial z_{1,\emptyset}=z_{1,\{0\}}z_{0,\emptyset}.
\]
Indeed, differentiating $z_{1,\{0\}}e_0$ introduces a minus sign,
so these identities give $\partial b_1^u=0$. The differential on
the coefficient algebra records the condition that the prescribed
expression be a cycle after the earlier null $e_0$ has been adjoined.
An actual evaluation may send some of these coefficients to zero.

\begin{lemma}\label{universal:triangular}
Definition~\ref{universal:coefficients} determines dg algebras
$R^a$, $K_-^a$, and $K^a$. The differential of each $z_{j,S}$ is
decomposable and uses only earlier generators when the generators
are ordered by increasing $j$ and, for a fixed $j$, by decreasing
binary order on subsets of $[j)$. There is a dg algebra evaluation
\begin{equation}\label{universal:evaluation}
 R^a\longrightarrow A|_\Lambda=h_\Lambda(\one)
\end{equation}
that sends $z_{j,S}$ to the coefficient of $e_S$ in $b_j$.
The constructions and evaluations extend those at every earlier stage.
\end{lemma}

\begin{proof}
Binary order compares two subsets at their largest differing index.
In a term of $\partial e_S$, one replaces an index $h\in S$ by a
subset of $[h)$. If the resulting exterior monomial is nonzero, its
index set is strictly smaller than $S$ in binary order. Therefore the
coefficient equation for $e_U$ in $\partial b_j^u$ expresses
$\partial z_{j,U}$ in terms of coefficients $z_{j,S}$ with $S>U$
and coefficients $z_{h,T}$ with $h<j$. The signs are fixed by
\[
 \partial(z_{j,S}e_S)
 = (\partial z_{j,S})e_S
    +(-1)^{|S|}z_{j,S}\partial e_S.
\]
Each nonzero summand in the expression for $\partial z_{j,U}$
contains one coefficient of each of these two kinds. It is thus
decomposable, with the required triangular ordering.

Inductively suppose that $\partial^2 e_h=0$ for $h<j$.
The square of an odd derivation is a derivation, since its mixed
terms cancel. Applying $\partial$ to $\partial b_j^u=0$ consequently
gives
\[
 0=\partial^2 b_j^u
   =\sum_{S\subseteq[j)}(\partial^2 z_{j,S})e_S.
\]
Independence of the exterior basis gives
$\partial^2z_{j,S}=0$ for every $S$, and the defining cycle equation
gives $\partial^2e_j=0$. This proves the assertion by induction.
The actual coefficients of $b_j$ satisfy precisely these equations,
so they define Equation~\eqref{universal:evaluation}; in particular
$z_{0,\emptyset}$ maps to $\varpi u_0$. All equations for an earlier
stage are retained when a new value of $j$ is added.
\end{proof}

Define an additive mass function on $\Lambda$ by
$\sigma(w_i)=3^i$. Every displayed generator has strictly positive
integral mass, since
\begin{equation}\label{universal:positive-mass}
 \sigma(\gamma_{j,S})
 \geq 3^j-\sum_{h<j}3^h=\frac{3^j+1}{2}>0,
 \qquad \sigma(w_j)=3^j>0.
\end{equation}
Give each generator $z_{j,S}$ and $e_j$ length one. We filter
$R^a$, $K_-^a$, and $K^a$ decreasingly by length: $\mathcal F_s$
is spanned by monomials of length at least $s$. At a fixed weight
$\gamma$, the filtration vanishes for $s>\sigma(\gamma)$ and is
constant for $s\leq0$. In particular it is bounded at each weight.

\begin{lemma}\label{universal:filtered-cells}
The length filtrations are dg filtrations. The associated graded
algebra $\gr R^a$ has zero differential, and in $\gr K^a$ one has
\[
 \partial e_j=z_{j,\emptyset}.
\]
As a filtered $R^a$-module, $K^a$ is built by finitely many free
cells indexed by the subsets $S\subseteq\{0,\ldots,a\}$, in
increasing binary order. The cell $e_S$ has chain degree $|S|$,
weight $\sum_{h\in S}w_h$, and filtration order $|S|$.
In particular, $K^a$ and $K_-^a$ are perfect $R^a$-modules.
\end{lemma}

\begin{proof}
The differential of a coefficient generator increases length,
by Lemma~\ref{universal:triangular}. The summand
$z_{j,S}e_S$ of $\partial e_j$ has length $1+|S|$, so only
$S=\emptyset$ contributes at the original length one.
For the module cell $e_S$, differentiating one factor $e_h$
and using the term indexed by $T\subseteq[h)$ gives total length
$|S|+|T|$. This is at least $|S|$, and the resulting exterior
basis vector is earlier in binary order. Thus the differential
is a filtered attachment to earlier module cells. There are
$2^{a+1}$ basis vectors, including the initial one indexed by
the empty set.
\end{proof}

\subsection{Filtered models at the primes}

The primewise models have two distinct roles. The coefficient algebras
$R_p^a$ will be actual commutative algebras, so their permutation powers
are defined. The objects $K_p^a$ will be finite-cell modules over them;
their role is to test vanishing at weights growing with $p$. They need
no primewise algebra structure. The rational Koszul algebra $K^a$
already has a commutative algebra structure, which will supply the
Amitsur construction in Section~\ref{sec:normalization}.

We retain the positive-mass filtration in both primewise models.
A filtered object is a diagram
$\mathcal F_s\to\mathcal F_{s-1}$ indexed by $s\in\Z$, with
convolution tensor product along addition. A step at order $s$
is constant at indices at most $s$ and zero at larger indices.
The tensor product of steps at orders $s$ and $t$ is a step at
order $s+t$.

We use filtrations supported in the weight monoid generated by
the relevant generator weights, constant below order zero,
and zero at orders greater than $\sigma(\gamma)$ in weight
$\gamma$. These conditions are preserved by colimits and tensor
products. For tensor products, one may decompose a bounded
filtration in each weight into finitely many extensions of
steps, starting at its highest nonzero order. On steps the
assertion follows from additivity of mass and of filtration order.

\begin{lemma}\label{universal:associated-graded}
On filtered spectra, the functor
\[
 (\gr\mathcal F)_s
   =\cofib(\mathcal F_{s+1}\longrightarrow\mathcal F_s)
\]
preserves colimits and is symmetric monoidal. It therefore preserves
free commutative algebras and algebra pushouts. The comparison
\[
 h_\Lambda(\gr X_p)\simeq\gr h_\Lambda(X_p)
\]
respects the algebra and module pairings whenever they are present.
\end{lemma}

\begin{proof}
The associated graded pairing is obtained by taking the cofiber
of the pushout-product of the filtration arrows. On steps this
gives the claimed tensor equivalence. Steps generate the filtered
category under colimits, and both constructions preserve colimits,
so the equivalence holds in general, with its symmetric monoidal
compatibilities. The assertions for free algebras and their
pushouts follow.

For the last assertion, fix a finite tuple of weights and
filtration orders. The tensor cube for the filtration arrows
maps to the arrow at the sum of those orders: every vertex
with at least one order increased maps into filtration at
least one higher than that sum. Taking the cube's
pushout-product and cofiber constructs the associated graded
pairing. The exact lax pairings defining $h_\Lambda$ commute
with these finite cofibers and produce the same construction.
This verifies the assertion for all finite tuples, together
with the compatibility under composition of pairings. It
does not require $h_\Lambda$ to commute with infinite
filtration colimits.
\end{proof}

\begin{proposition}\label{universal:primewise-models}
For the chosen cycles through $b_a$, there exist the following data
as germs along $\U$:
\begin{enumerate}
\item Filtered $\Lambda$-graded commutative ring spectra $R_p^a$,
with an equivalence
$R^a\simeq h_\Lambda(R_p^a)$ at every filtration order, compatible
with the filtered algebra structures.
\item Maps of whole unfiltered graded commutative algebras
$R_p^a\to S_{(p)}$, where the sphere is constant in the weight
grading, together with a homotopy identifying their normalized
evaluation with Equation~\eqref{universal:evaluation}.
\item Filtered finite-cell $R_p^a$-modules $K_p^a$, with an
equivalence $K^a\simeq h_\Lambda(K_p^a)$ of filtered modules.
\end{enumerate}
The algebra models, their comparisons, and the evaluation
homotopies may be chosen relative to the preceding stage.
Thus the maps $R^{a-1}\to R^a$ and $R_p^{a-1}\to R_p^a$
and their evaluations form compatible diagrams of algebras.
No algebra structure is required on $K_p^a$.
\end{proposition}

\begin{proof}
We first construct the algebras in the triangular generator order
of Lemma~\ref{universal:triangular}. Start at the unit, in weight
and filtration order zero. Suppose that the preceding generators
have been constructed, and let $z=z_{j,S}$, $r=|S|$, and
$\gamma=\gamma_{j,S}$. Its boundary $\partial z$ is a cycle of
chain degree $-r-1$, weight $\gamma$, and filtration at least two.
The comparison already established at this fixed weight and
filtration lifts the boundary to maps
\begin{equation}\label{universal:attaching-map}
 S^{d(\gamma)-r-1}
    \longrightarrow(\mathcal F_2R_p^{\mathrm{old}})(\gamma),
\end{equation}
together with a path matching the normalized class with
$\partial z$. Compose with the map to $\mathcal F_1$ and attach
a free commutative algebra cell that nulls this map. Explicitly,
push out the free algebra on the indicated sphere, placed as
a step at order one, along its augmentation to the unit.
The support and filtration bounds are preserved.

After taking associated graded, the attaching map in order one
is null by its specified factorization through $\mathcal F_2$.
Lemma~\ref{universal:associated-graded} identifies the result
with free adjunction of the sphere
$S^{d(\gamma)-r}$ in weight $\gamma$ and order one. Repeating
this construction gives a primewise free algebra on all the
generator spheres as the associated graded of $R_p^a$.

We verify carefully that the comparison with the filtered dg
algebra extends at each step. Over the $\Q$-algebra $D$, free
algebras on free step generators are computed by ordinary
graded symmetric powers, with their step orders added. For
a cycle generator $x$, the free disk algebra with a generator
$z$ satisfying $\partial z=x$, both placed at order one,
resolves the filtered unit over the free cycle algebra.
Indeed every component of positive word length is
contractible. If the null generator is even, the needed
contraction divides by its positive multiplicity, which is
invertible in $D$. Moreover, filtering the disk algebra by
the number of null generators gives successive free step
modules over the cycle algebra. The disk therefore computes
the derived relative tensor, and its base change is exactly
the extension by $z$ with the specified differential and
length filtration.

These computations use derived tensor products. To see
directly that the strict monomial models compute them,
order monomials at each weight by their highest differing
generator multiplicity in the triangular list. There are
finitely many such monomials by positive mass, and their
differentials use earlier ones. For $K^a$, place the $e_j$
after all coefficient generators, in increasing $j$.
The resulting filtered complexes are built from free
$D$-module steps. Their strict tensor pairings thus
represent the derived pairings as well.

The old comparison and the tautological null in the spectral
attachment now induce a map from the dg pushout to the
normalized spectral pushout. We use the chosen matching
path from Equation~\eqref{universal:attaching-map} inside
$\mathcal F_2$. In order one modulo order two, the new
generator is the difference of two nulls: the attached null
and the null obtained by projecting the boundary from
$\mathcal F_2$. The matching path respects the latter null.
Consequently the comparison sends the new associated graded
generator to the spectral generator supplied by the attachment.
Taking cofibers and projecting the filtration commute with
normalized Hom, so this is also an identification of the
chosen generator classes, not merely of the objects.

At any fixed weight, the associated graded spectral algebra
is a finite sum of extended powers of the generator spheres,
with arities bounded by the mass of that weight independently
of $p$. For sufficiently large $p$, the symmetric groups
in these terms have invertible order. An even sphere
repeated $k$ times contributes a sphere of the summed degree:
the fiber inclusion into its homotopy orbits is an equivalence
after $p$-localization. An odd sphere repeated at least twice
contributes zero. To justify both assertions, its extended
power is the Thom spectrum of the permutation sphere over
$B\Sigma_k$, and its ordinary homology is the group homology
with the permutation orientation character. When
$p\nmid k!$, higher group homology vanishes. The degree-zero
coinvariants are $\Z_{(p)}$ in the even case and zero in the
odd case, since a transposition acts by $-1$ and $2$ is a
unit. The spectra are bounded below after translating the
sphere degree, so the homology statements imply the claimed
equivalences by stable Hurewicz.

Products of distinct generator factors are treated separately.
Multiplication of their generator classes uses the fiber
inclusions. The coherent sphere grading and
$\map_{\cC}(\one,\one)=HD$ identify their normalized Hom
with the expected free graded-commutative basis over $D$.
The extra chain shifts are fixed at each fixed weight, so
the fixed-stem calculation applies. Our comparison is
therefore an equivalence on associated gradeds, and hence
on the bounded filtration at every weight. This completes
the induction constructing the filtered algebra comparison.

It remains to realize the evaluation, including its specified
nulls. If $(X_p)$ is any sequence of graded commutative
algebras, there is a comparison
\begin{equation}\label{universal:algebra-mapping}
 \prod_{\U}\Map_{\CAlg^\Lambda}(R_p^a,X_p)
 \longrightarrow
 \Map_{\CAlg_D^\Lambda}(R^a,h_\Lambda(X_p)),
\end{equation}
where the source uses maps of whole graded systems, after
forgetting filtration. This is an equivalence by induction
over the finite list of free algebra cells. For one cell,
extending a map means choosing a null of the image of its
attaching map. Its mapping space is the corresponding
homotopy pullback of mapping spaces from the generator
sphere at its specified weight. Those sphere mapping
spaces agree by the definition of $h_\Lambda$. The initial
unit case agrees as well, and finite homotopy pullbacks
commute with the ultraproduct. The comparison constructed
above sends attached-null data to the attached-null data
in these pullbacks, since it used both the tautological
null and the specified matching path.

Forgetting filtration in this argument does preserve the
free-cell pushouts: it is colimit towards decreasing
filtration order, a colimit-preserving symmetric monoidal
functor, as one sees on steps. Here it takes the constant
value at orders at most zero. Apply
Equation~\eqref{universal:algebra-mapping} to the constant
graded sphere target and the evaluation in
Equation~\eqref{universal:evaluation}. It supplies the
required primewise evaluations and the matching homotopy.
The same comparison commutes with restriction to earlier
cells. It is therefore an equivalence on the homotopy fibers
over already chosen maps, including their comparison paths.
The new evaluation can thus extend the preceding evaluation
with its chosen compatibility, which proves the assertion
about diagrams across stages.

Finally, construct $K_p^a$ using the finite filtered module
cells from Lemma~\ref{universal:filtered-cells}. A free module
cell over $R_p^a$ with generator of weight $\gamma$, order
$s$, and ordinary degree $d(\gamma)+u$, for fixed $u\in\Z$,
normalizes to the corresponding free module shift over
$h_\Lambda(R_p^a)$ with chain degree $u$. This uses the
coherent sphere grading. A boundary is, by adjunction, a
map from one sphere into the specified weight and
filtration of the preceding module. Lift that map and
choose a path matching the boundary. Taking its cofiber
extends the module comparison. Exactness proves the
comparison at each stage and gives the compatible
associated graded comparison.
\end{proof}

\subsection{A comparison of whole primewise modules}

The fixed-weight equivalences just proved do not, by themselves,
control a weight that depends on $p$. We now obtain a stronger
statement using the finite number of module cells. This is the
point at which the whole graded systems in
Proposition~\ref{universal:primewise-models} are essential.

\begin{proposition}\label{universal:whole-module}
On one $\U$-large set of primes there is an equivalence of whole
$(\text{weight},\text{order})$-graded $\gr R_p^a$-modules
\begin{equation}\label{eq:whole-module}
 \gr K_p^a\simeq
 \bigotimes_{j\leq a}^{\gr R_p^a}
 \cofib\bigl(\text{multiplication by }z_{j,\emptyset}
                         \text{ on }\gr R_p^a\bigr).
\end{equation}
Each multiplication map has source the free module whose
generator has weight $w_j$, order one, and ordinary degree
$d_j$. On the right, $\gr R_p^a$ is the free graded algebra
of all the generator spheres, with the splittings constructed
in Proposition~\ref{universal:primewise-models}.
\end{proposition}

\begin{proof}
After applying $h_\Lambda$, the right side becomes the tensor
of two-term cofibers for the classes $z_{j,\emptyset}$ in
$\gr R^a$. The tensor comparison is an equivalence for
these finite free module cells: it is the free-shift
identification for one cell and extends by finite cofibers.
This is the Koszul module with differential
$\partial e_j=z_{j,\emptyset}$, hence agrees with the
normalized left side by
Lemma~\ref{universal:filtered-cells}. A change in a sphere
identification can multiply a chosen basis element by a
unit and does not change this equivalence.

To lift that equivalence, consider the ultraproduct of
categories of whole $(\text{weight},\text{order})$-graded
modules over $\gr R_p^a$. There is a natural comparison
from its mapping spectra to mapping spectra of graded
modules over $h_\Lambda(\gr R_p^a)$. For a free source
at one fixed weight, order, and extra integer suspension,
both mapping spectra are exactly the same normalized
evaluation of the target at that weight and order.
Induction through cofiber sequences proves full faithfulness
for sources built from finitely many such free shifts,
against any whole-system target.

Both modules in Equation~\eqref{eq:whole-module} satisfy
this finiteness condition. For the left side, take
associated graded of the finite filtered cell construction.
For the right side, expand the tensor of the finitely many
two-term cofibers. Their formal weight and order shifts and
their extra integer suspensions are fixed independently
of $p$. Full faithfulness lifts the normalized equivalence,
its inverse, and the two homotopies identifying their
composites with the respective identities. These are
finitely many maps and homotopies of whole systems. After
restricting to one $\U$-large prime set, they are inverse
equivalences at each prime as whole graded module maps.
Thus the same prime set works for all weights; no
intersection over the lattice of weights is taken.
\end{proof}

Equation~\eqref{eq:whole-module} lets us examine the same modules
at $p$-dependent weights. The next section estimates these
weights by analyzing the extended powers of each generator
sphere and the multiplication cofibers of the empty-subset
generators.
